%==============================================================================
\chapter{Test plan and validation}
\label{chap:tests}
%==============================================================================
The test strategy covers the business functions, the management constraints and
the technical requirements. Each test case specifies the target function, the
preconditions, the steps and the expected result.

\section{Test cases}
\begin{xltabular}{\textwidth}{>{\bfseries}p{1.4cm} p{3.6cm} p{4.2cm} X}
\toprule
\textbf{ID} & \textbf{Function} & \textbf{Steps} & \textbf{Expected result} \\
\midrule
\endhead
TC01 & Hive CRUD & Create, read, modify, delete a hive & Compliant and persisted operations \\
TC02 & Frame constraint & Add an 11th frame to a body & Rejected, maximum of 10 frames \\
TC03 & Super constraint & Add a 6th super & Rejected, maximum of 5 supers \\
TC04 & Base constraint & Remove the last frame of the body & Rejected, at least 1 frame \\
TC05 & Visit schedule & Propose then approve a visit & Approved status recorded \\
TC06 & Visit report & Fill in all fields and save & Report persisted, cell carried out \\
TC07 & Production alert & Exceed the production threshold & Harvest alert displayed \\
TC08 & Drop alert & Simulate a population drop & Inspection alert displayed \\
TC09 & ROI calculation & Enter costs and production & ROI computed correctly \\
TC10 & Honey API & Call getZummHoneyActualQuantity & Correct quantity returned \\
TC11 & Internationalisation & Switch FR, EN then AR & Texts translated via the XML file \\
TC12 & Configuration & Modify a threshold in ConfigZumm.ini & New threshold taken into account \\
TC13 & Authentication & Log in via Google OAuth & Access granted \\
TC14 & Security & Verify SSL and X.509 certificate & Encrypted and validated exchanges \\
TC15 & Export & Export the records & CSV and TXT files generated \\
TC16 & Offline mode & Enter without a network then reconnect & Successful synchronisation \\
\bottomrule
\end{xltabular}

\section{Traceability matrix}
The matrix links each requirement to its use case, to the classes and tables
that carry it, to the screen (mockup) concerned and to the test cases that
validate it. It ensures end-to-end traceability, from the need to the
verification.
{\small
\begin{xltabular}{\textwidth}{>{\bfseries}X C{1.3cm} p{3.3cm} p{2.5cm} p{2.6cm}}
\toprule
\textbf{Requirement} & \textbf{UC} & \textbf{Classes / Tables} & \textbf{Screen} & \textbf{Tests} \\
\midrule
\endhead
Entity management and constraints & --- & Hive, Compartment, Frame & Hive form (fig.~\ref{fig:wf-ruche}) & TC01--TC04 \\
Planning and approval & UC1 & Schedule, Agent & Calendar (fig.~\ref{fig:wf-calendrier}) & TC05 \\
Visit execution and report & UC2 & Visit, Photo & Report (fig.~\ref{fig:wf-rapport}) & TC06 \\
Production dashboard & UC3 & Harvest, Measurement, Threshold & Prod. dashboard (fig.~\ref{fig:wf-production}) & TC07, TC09 \\
Health / drop alert & UC4 & Measurement, Threshold & Calendar & TC08 \\
API web service (honey) & UC5 & Hive, Harvest & --- & TC10 \\
Internationalisation & --- & (i18n XML) & All & TC11 \\
Threshold configuration & --- & Threshold & Prod. dashboard & TC12 \\
OAuth authentication & --- & Agent & Login (fig.~\ref{fig:seq-oauth}) & TC13 \\
Exchange security & --- & (SSL / X.509) & --- & TC14 \\
Portability and export & --- & (CSV/TXT export) & All & TC15 \\
Offline mode & UC2 & Visit (local queue) & Report & TC16 \\
\bottomrule
\end{xltabular}}
